% Options for packages loaded elsewhere
\PassOptionsToPackage{unicode}{hyperref}
\PassOptionsToPackage{hyphens}{url}
\documentclass[
]{article}
\usepackage{xcolor}
\usepackage{amsmath,amssymb}
\setcounter{secnumdepth}{-\maxdimen} % remove section numbering
\usepackage{iftex}
\ifPDFTeX
  \usepackage[T1]{fontenc}
  \usepackage[utf8]{inputenc}
  \usepackage{textcomp} % provide euro and other symbols
\else % if luatex or xetex
  \usepackage{unicode-math} % this also loads fontspec
  \defaultfontfeatures{Scale=MatchLowercase}
  \defaultfontfeatures[\rmfamily]{Ligatures=TeX,Scale=1}
\fi
\usepackage{lmodern}
\ifPDFTeX\else
  % xetex/luatex font selection
\fi
% Use upquote if available, for straight quotes in verbatim environments
\IfFileExists{upquote.sty}{\usepackage{upquote}}{}
\IfFileExists{microtype.sty}{% use microtype if available
  \usepackage[]{microtype}
  \UseMicrotypeSet[protrusion]{basicmath} % disable protrusion for tt fonts
}{}
\makeatletter
\@ifundefined{KOMAClassName}{% if non-KOMA class
  \IfFileExists{parskip.sty}{%
    \usepackage{parskip}
  }{% else
    \setlength{\parindent}{0pt}
    \setlength{\parskip}{6pt plus 2pt minus 1pt}}
}{% if KOMA class
  \KOMAoptions{parskip=half}}
\makeatother
\usepackage{color}
\usepackage{fancyvrb}
\newcommand{\VerbBar}{|}
\newcommand{\VERB}{\Verb[commandchars=\\\{\}]}
\DefineVerbatimEnvironment{Highlighting}{Verbatim}{commandchars=\\\{\}}
% Add ',fontsize=\small' for more characters per line
\newenvironment{Shaded}{}{}
\newcommand{\AlertTok}[1]{\textcolor[rgb]{1.00,0.00,0.00}{\textbf{#1}}}
\newcommand{\AnnotationTok}[1]{\textcolor[rgb]{0.38,0.63,0.69}{\textbf{\textit{#1}}}}
\newcommand{\AttributeTok}[1]{\textcolor[rgb]{0.49,0.56,0.16}{#1}}
\newcommand{\BaseNTok}[1]{\textcolor[rgb]{0.25,0.63,0.44}{#1}}
\newcommand{\BuiltInTok}[1]{\textcolor[rgb]{0.00,0.50,0.00}{#1}}
\newcommand{\CharTok}[1]{\textcolor[rgb]{0.25,0.44,0.63}{#1}}
\newcommand{\CommentTok}[1]{\textcolor[rgb]{0.38,0.63,0.69}{\textit{#1}}}
\newcommand{\CommentVarTok}[1]{\textcolor[rgb]{0.38,0.63,0.69}{\textbf{\textit{#1}}}}
\newcommand{\ConstantTok}[1]{\textcolor[rgb]{0.53,0.00,0.00}{#1}}
\newcommand{\ControlFlowTok}[1]{\textcolor[rgb]{0.00,0.44,0.13}{\textbf{#1}}}
\newcommand{\DataTypeTok}[1]{\textcolor[rgb]{0.56,0.13,0.00}{#1}}
\newcommand{\DecValTok}[1]{\textcolor[rgb]{0.25,0.63,0.44}{#1}}
\newcommand{\DocumentationTok}[1]{\textcolor[rgb]{0.73,0.13,0.13}{\textit{#1}}}
\newcommand{\ErrorTok}[1]{\textcolor[rgb]{1.00,0.00,0.00}{\textbf{#1}}}
\newcommand{\ExtensionTok}[1]{#1}
\newcommand{\FloatTok}[1]{\textcolor[rgb]{0.25,0.63,0.44}{#1}}
\newcommand{\FunctionTok}[1]{\textcolor[rgb]{0.02,0.16,0.49}{#1}}
\newcommand{\ImportTok}[1]{\textcolor[rgb]{0.00,0.50,0.00}{\textbf{#1}}}
\newcommand{\InformationTok}[1]{\textcolor[rgb]{0.38,0.63,0.69}{\textbf{\textit{#1}}}}
\newcommand{\KeywordTok}[1]{\textcolor[rgb]{0.00,0.44,0.13}{\textbf{#1}}}
\newcommand{\NormalTok}[1]{#1}
\newcommand{\OperatorTok}[1]{\textcolor[rgb]{0.40,0.40,0.40}{#1}}
\newcommand{\OtherTok}[1]{\textcolor[rgb]{0.00,0.44,0.13}{#1}}
\newcommand{\PreprocessorTok}[1]{\textcolor[rgb]{0.74,0.48,0.00}{#1}}
\newcommand{\RegionMarkerTok}[1]{#1}
\newcommand{\SpecialCharTok}[1]{\textcolor[rgb]{0.25,0.44,0.63}{#1}}
\newcommand{\SpecialStringTok}[1]{\textcolor[rgb]{0.73,0.40,0.53}{#1}}
\newcommand{\StringTok}[1]{\textcolor[rgb]{0.25,0.44,0.63}{#1}}
\newcommand{\VariableTok}[1]{\textcolor[rgb]{0.10,0.09,0.49}{#1}}
\newcommand{\VerbatimStringTok}[1]{\textcolor[rgb]{0.25,0.44,0.63}{#1}}
\newcommand{\WarningTok}[1]{\textcolor[rgb]{0.38,0.63,0.69}{\textbf{\textit{#1}}}}
\usepackage{longtable,booktabs,array}
\newcounter{none} % for unnumbered tables
\usepackage{calc} % for calculating minipage widths
% Correct order of tables after \paragraph or \subparagraph
\usepackage{etoolbox}
\makeatletter
\patchcmd\longtable{\par}{\if@noskipsec\mbox{}\fi\par}{}{}
\makeatother
% Allow footnotes in longtable head/foot
\IfFileExists{footnotehyper.sty}{\usepackage{footnotehyper}}{\usepackage{footnote}}
\makesavenoteenv{longtable}
\setlength{\emergencystretch}{3em} % prevent overfull lines
\providecommand{\tightlist}{%
  \setlength{\itemsep}{0pt}\setlength{\parskip}{0pt}}
\usepackage{bookmark}
\IfFileExists{xurl.sty}{\usepackage{xurl}}{} % add URL line breaks if available
\urlstyle{same}
\hypersetup{
  pdftitle={Factorial Barrier Closure: 26! = (1)\_\{26\} from 26-fold Radial Derivative (G8 Closure)},
  pdfauthor={Daniel T. Murphy},
  hidelinks,
  pdfcreator={LaTeX via pandoc}}

\title{Factorial Barrier Closure: 26! = (1)\_\{26\} from 26-fold Radial
Derivative (G8 Closure)}
\author{Daniel T. Murphy}
\date{2026-05-10}

\begin{document}
\maketitle

\section{\texorpdfstring{PAPER\_1161 -- Factorial Barrier Closure:
\((26!)^2\) from 26-fold Radial
Derivative}{PAPER\_1161 -- Factorial Barrier Closure: (26!)\^{}2 from 26-fold Radial Derivative}}\label{paper_1161-factorial-barrier-closure-262-from-26-fold-radial-derivative}

\textbf{Author:} Daniel T. Murphy (daniel.murphy00@gmail.com)
\textbf{Framework:} UQFF v5.78 -- Star-Magic Physics \textbf{Source:}
Direct identification from existing codebase
(\href{QCalcGeom.cpp\#L431}{QCalcGeom.cpp L431-458};
\href{dpm_vacuum_manifold.py\#L1251}{dpm\_vacuum\_manifold.py
L1251-1258}). \textbf{Integration Date:} May 10, 2026 (Session 248)
\textbf{Classification:} Lagrangian Gap Closure -- G8 of
\texttt{\_lagrangian\_rederivation\_outline.py}

\begin{center}\rule{0.5\linewidth}{0.5pt}\end{center}

\[\boxed{\;26! \;=\; (1)_{26} \;=\; \frac{d^{26}}{dr^{26}}\!\left(\frac{1}{r}\right) \cdot (-1)^{26}\,r^{27} \;=\; \prod_{k=1}^{26} k \;=\; 4.0329 \times 10^{26}\;}\]

The \((26!)^2\) factor in the \(G\) closure (PAPER\_593) is the
\textbf{square} of this single Pochhammer symbol -- one factor from the
kinetic-term variation of the BSFG action, one from the source coupling,
both iterated to the 26th radial order at the bosonic critical dimension
\(D_{\rm crit}=26\).

\begin{center}\rule{0.5\linewidth}{0.5pt}\end{center}

\subsection{Abstract}\label{abstract}

We close gap \textbf{G8} of the Lagrangian re-derivation outline by
identifying the \(26!\) factorial barrier in the \(G\) closure
(PAPER\_593) as the \textbf{Pochhammer rising factorial \((1)_{26}\)}
arising from the 26-fold iterated radial derivative of the Newtonian
Green's function \(1/r\) at the bosonic critical dimension
\(D_{\rm crit} = 26\). The identification is \textbf{already explicit in
the codebase} (\href{QCalcGeom.cpp\#L450}{QCalcGeom.cpp L450},
\href{QCalcGeom.cpp\#L458}{L458}) but was not previously catalogued as
the structural origin of the \(26!\) appearing in \(G_{\rm UQFF}\). The
square \((26!)^2\) in the \(G\) closure denominator is the natural
consequence of varying a 26th-order multipole coupling on \textbf{both}
sides of the Einstein-like field equation: kinetic term and source term
each contribute one factor.

\begin{center}\rule{0.5\linewidth}{0.5pt}\end{center}

\subsection{1. The Gap}\label{the-gap}

The Newton-constant closure (PAPER\_593, Session 240) takes the form:
\[G_{\rm UQFF} \;=\; \frac{2\pi \cdot 26^3 \cdot \Phi_{\rm res}}{[\mathrm{SSq}]^3 \cdot (26!)^2} \cdot \frac{v_F^5}{E_0 \cdot f_{\rm THz}}\]

The factor \((26!)^2 \approx 1.626 \times 10^{53}\) supplies the
\(\sim 10^{53}\) hierarchy gap between the Planck-scale resonance
density and the observed gravitational coupling
(\href{AXIOMS_AND_THEOREMS.md\#L288}{AXIOMS\_AND\_THEOREMS.md L288}).
Until Session 248 this factorial barrier was asserted
phenomenologically, not derived. It was the eighth and final identified
gap in \texttt{\_lagrangian\_rederivation\_outline.py}.

\subsection{2. The Identification (already in the
codebase)}\label{the-identification-already-in-the-codebase}

Two existing codebase entries already perform the identification but
were never explicitly catalogued as the G8 closure:

\subsubsection{\texorpdfstring{2.1
\href{QCalcGeom.cpp\#L431}{QCalcGeom.cpp L431} -- Pochhammer rising
factorial}{2.1 QCalcGeom.cpp L431 -- Pochhammer rising factorial}}\label{qcalcgeom.cpp-l431-pochhammer-rising-factorial}

\begin{Shaded}
\begin{Highlighting}[]
\CommentTok{// Pochhammer rising factorial: (k)\_\{26\} = k*(k+1)*...*(k+25) = (k+25)!/(k{-}1)!}
\end{Highlighting}
\end{Shaded}

For \(k = 1\): \[(1)_{26} \;=\; 1 \cdot 2 \cdot 3 \cdots 26 \;=\; 26!\]
The Pochhammer symbol arises in the \textbf{iterated raising operator}
of the BSFG ladder representation -- each rung of the 26D atlas
contributes one factor.

\subsubsection{\texorpdfstring{2.2
\href{QCalcGeom.cpp\#L450}{QCalcGeom.cpp L450} -- 26th radial
derivative}{2.2 QCalcGeom.cpp L450 -- 26th radial derivative}}\label{qcalcgeom.cpp-l450-26th-radial-derivative}

\begin{Shaded}
\begin{Highlighting}[]
\CommentTok{// U\_g approx G*M\_sun/r  =\textgreater{}  d\^{}\{26\}/dr\^{}\{26\}(G*M\_sun/r) = 26! * G*M\_sun / r\^{}\{27\}}
\end{Highlighting}
\end{Shaded}

The 26th derivative of the Newtonian potential picks up exactly \(26!\)
as the Leibniz coefficient. This is the \textbf{leading multipole term}
of the BSFG generalised gravitational action when expanded to 26 radial
moments (one moment per critical dimension).

\subsubsection{\texorpdfstring{2.3
\href{dpm_vacuum_manifold.py\#L1251}{dpm\_vacuum\_manifold.py
L1251-1258} -- 26D geometric
folding}{2.3 dpm\_vacuum\_manifold.py L1251-1258 -- 26D geometric folding}}\label{dpm_vacuum_manifold.py-l1251-1258-26d-geometric-folding}

\begin{Shaded}
\begin{Highlighting}[]
\CommentTok{"""UQFF 26D geometric folding: crossing radius scaled by inverse 26th{-}root of 26!.}
\CommentTok{F26(r) = r * (26!)\^{}({-}1/13) * S26\_3 * Phi\_resonance}
\CommentTok{(26!)\^{}({-}1/13) provides the characteristic sub{-}Planck geometric scale."""}
\end{Highlighting}
\end{Shaded}

The exponent \(1/13 = 2/26\) confirms that the \(26!\) enters via a
\textbf{square} of the per-rung product when projected onto the 26-rung
ladder (\(1/13\) on each radial axis after the 13-pair factorisation).

\subsection{\texorpdfstring{3. The Square: \((26!)^2\) from Two-Sided
Variation}{3. The Square: (26!)\^{}2 from Two-Sided Variation}}\label{the-square-262-from-two-sided-variation}

The Einstein-like field equation in BSFG form, schematically,
\[\hat{\mathcal{D}}^{(26)}_{\rm kin} \cdot G \;=\; \hat{\mathcal{D}}^{(26)}_{\rm src} \cdot T,\]
varies the action \textbf{twice} under iterated 26th-order radial
derivatives:

{\def\LTcaptype{none} % do not increment counter
\begin{longtable}[]{@{}llr@{}}
\toprule\noalign{}
Side & Operator & Pochhammer factor \\
\midrule\noalign{}
\endhead
\bottomrule\noalign{}
\endlastfoot
Kinetic (LHS) & \(\partial_r^{26}\,(1/r)\) & \(26!\) \\
Source (RHS) & \(\partial_r^{26}\,T(r)\) & \(26!\) \\
\end{longtable}
}

Solving for \(G\):
\[G \;=\; \frac{(\text{numerator})}{(\text{kinetic Pochhammer}) \cdot (\text{source Pochhammer})} \;=\; \frac{\dots}{26! \cdot 26!} \;=\; \frac{\dots}{(26!)^2}\]

This is the structural origin of the \textbf{square}, not an additional
free input. Once the bosonic critical dimension \(D_{\rm crit}=26\) is
fixed (a textbook string-theory result), the square is automatic from
the standard variational principle.

\subsection{\texorpdfstring{4. Why \(D_{\rm crit} = 26\)? (Already
textbook)}{4. Why D\_\{\textbackslash rm crit\} = 26? (Already textbook)}}\label{why-d_rm-crit-26-already-textbook}

The bosonic critical dimension \(D_{\rm crit} = 26\) is the
\textbf{only} spacetime dimension at which the Polyakov action yields a
ghost-free, conformally invariant string spectrum. This is a classical
result of Kaku, Polchinski, GSW; it requires zero free parameters in
UQFF and is referenced explicitly in
\href{CondensedPhysics2.py\#L28610}{CondensedPhysics2.py L28610} and
\href{COMPLETE_UQFF_EQUATIONS_REFERENCE.md\#L1065}{COMPLETE\_UQFF\_EQUATIONS\_REFERENCE.md
L1065}: \emph{``Critical dimension: \(D = 26\) (= VDS 26-rung ladder,
\(\mathrm{Li}_{26}([SSq])\)).''}

The chain is:
\[D_{\rm crit}=26 \;\xrightarrow{\;\text{compactify}\;}\; D_{\rm super}=10 \;\xrightarrow{\;\text{compactify}\;}\; D_{\rm BSFG}=6 \;\xrightarrow{\;\text{compactify}\;}\; D_{\rm phys}=4\]

PAPER\_1159 used \(D_{\rm BSFG}=6\) for \(\Phi_{\rm res}\). PAPER\_1160
used \(D_{\rm BSFG}-1=5\) for \(F_{\rm TRZ}\). \textbf{PAPER\_1161 uses
\(D_{\rm crit}=26\)} for the factorial barrier. All three are consistent
with the same \(26 \to 10 \to 6 \to 4\) flow.

\subsection{5. Numerical Match -- Trivially
Exact}\label{numerical-match-trivially-exact}

{\def\LTcaptype{none} % do not increment counter
\begin{longtable}[]{@{}
  >{\raggedright\arraybackslash}p{(\linewidth - 4\tabcolsep) * \real{0.3000}}
  >{\raggedleft\arraybackslash}p{(\linewidth - 4\tabcolsep) * \real{0.4000}}
  >{\raggedright\arraybackslash}p{(\linewidth - 4\tabcolsep) * \real{0.3000}}@{}}
\toprule\noalign{}
\begin{minipage}[b]{\linewidth}\raggedright
Quantity
\end{minipage} & \begin{minipage}[b]{\linewidth}\raggedleft
Value
\end{minipage} & \begin{minipage}[b]{\linewidth}\raggedright
Status
\end{minipage} \\
\midrule\noalign{}
\endhead
\bottomrule\noalign{}
\endlastfoot
\((1)_{26}\) Pochhammer &
\(403{,}291{,}461{,}126{,}605{,}635{,}584{,}000{,}000\) & exact \\
\(26!\) direct & \(4.0329 \times 10^{26}\) & exact \\
\(\partial_r^{26}(1/r)\cdot r^{27}\) at \(r=1\) & \(26!\) & exact
(Leibniz) \\
Factorial barrier \((26!)^2\) in \(G\) closure &
\(1.626 \times 10^{53}\) & exact \\
\end{longtable}
}

There is \textbf{no residual}: \(26!\) is an integer that arises exactly
from the standard 26-fold derivative formula. The \(G\) closure residual
(\(\sim 0.87\%\) with structural \(\Phi_{\rm res}=5/6\), see
\href{whitepapers/PAPER_1159_UQFF_Phi_Res_Codimension_Closure.md}{PAPER\_1159
§6}) is \textbf{entirely} attributable to the one-loop
\(\Phi_{\rm res}\) correction, not to the factorial barrier.

\subsection{\texorpdfstring{6. Cross-Check: Why Not
\(22!\)?}{6. Cross-Check: Why Not 22!?}}\label{cross-check-why-not-22}

The Lagrangian outline
(\href{_lagrangian_rederivation_outline.py\#L154}{L154-156})
hypothesised ``a combinatorial count of the 22-torus winding sectors,''
i.e., \(22!\). Numerical test:
\[22! \;=\; 1.1240 \times 10^{21}, \qquad \frac{22!}{26!} \;=\; 2.79 \times 10^{-6}\]
which would shift \(G\) by a factor of \(\sim 10^{11}\) --
\textbf{catastrophically wrong}. The factorial must come from the
\textbf{full critical dimension \(D_{\rm crit}=26\)}, not from the
compactified 22-torus. The compactification reduces the
\textbf{effective} dimensions (\(26 \to 10 \to 6 \to 4\)) but the action
variation retains the critical-dimension multipole order because the
underlying string worldsheet still has \(D_{\rm crit}=26\) degrees of
freedom.

This matches the \(D=26\) used in the VDS series
\(\mathrm{Li}_{26}([\mathrm{SSq}])\) and in the BSFG 26-rung ladder.

\subsection{7. Updated Catalog Status}\label{updated-catalog-status}

{\def\LTcaptype{none} % do not increment counter
\begin{longtable}[]{@{}lll@{}}
\toprule\noalign{}
Lagrangian gap & Status (pre-248) & Status (post-248) \\
\midrule\noalign{}
\endhead
\bottomrule\noalign{}
\endlastfoot
G1 (V(UA) polynomial) & OPEN & OPEN \\
G2 (beta\_i index) & OPEN & OPEN \\
G3 (DPM SO(2)) & OPEN & OPEN \\
G4 (T\^{}22 moduli) & OPEN & OPEN \\
G5 (KK tower) & OPEN & OPEN \\
G6 (Phi\_res ID) & CLOSED (PAPER\_1159) & CLOSED \\
G7 (F\_TRZ ID) & CLOSED (PAPER\_1160) & CLOSED \\
\textbf{G8 (26! emergence)} & \textbf{OPEN} & \textbf{CLOSED (this
paper, exact)} \\
\end{longtable}
}

\begin{quote}
\textbf{Session 253 Update (PAPER\_1166):} ALL 8 Lagrangian gaps now
closed. The \(26 \to 10 \to 6 \to 4\) critical-dimension flow noted
below proved comprehensive -- every subsequent closure (G3, G4, G5, G2,
G1) reduced to it.
\end{quote}

\textbf{Result:} \textbf{3 of 8 gaps closed}; 5 remain. All three
closures (\(\Phi_{\rm res}, F_{\rm TRZ}, 26!\)) trace back to the
\(26 \to 10 \to 6 \to 4\) critical-dimension flow. \textbf{Two integers
(\(D_{\rm crit}=26\) and \(D_{\rm BSFG}=6\)) now suffice} to
parameterise four formerly-free numerical inputs (\(\Phi_{\rm res}\),
\(F_{\rm TRZ}\), \(26!\), and the \(26^3\) prefactor in \(G\)).

\subsection{8. Conclusions}\label{conclusions}

\begin{itemize}
\tightlist
\item
  \(26!\) in the \(G\) closure is the Pochhammer \((1)_{26}\) from the
  26-fold iterated radial derivative of the Newtonian Green's function
  -- a standard Leibniz-rule calculation already encoded in
  \href{QCalcGeom.cpp\#L450}{QCalcGeom.cpp L450}.
\item
  The square \((26!)^2\) in \(G_{\rm UQFF}\) is the inevitable
  consequence of varying a 26th-order multipole on \textbf{both} the
  kinetic and source sides of the field equation.
\item
  The integer \(26\) comes from the \textbf{bosonic critical dimension}
  \(D_{\rm crit}=26\), a textbook string-theory result requiring zero
  free parameters.
\item
  Alternative candidates (\(22!\) from torus windings) fail by 11 orders
  of magnitude. The full \(D_{\rm crit}=26\) is required.
\item
  G8 closed; \textbf{3 of 8 Lagrangian gaps now closed}. Only G1-G5
  remain (V(UA) polynomial, beta\_i index, DPM SO(2), T\^{}22 moduli, KK
  tower).
\item
  Cumulative impact across PAPERs 1159-1161: 4 free numerical inputs
  reduced to 2 textbook integers (\(D_{\rm crit}=26\),
  \(D_{\rm BSFG}=6\)).
\end{itemize}

\begin{center}\rule{0.5\linewidth}{0.5pt}\end{center}

\subsection{§SM Anchors (G8 Compliance
Table)}\label{sm-anchors-g8-compliance-table}

{\def\LTcaptype{none} % do not increment counter
\begin{longtable}[]{@{}
  >{\raggedright\arraybackslash}p{(\linewidth - 8\tabcolsep) * \real{0.1875}}
  >{\raggedright\arraybackslash}p{(\linewidth - 8\tabcolsep) * \real{0.1875}}
  >{\raggedleft\arraybackslash}p{(\linewidth - 8\tabcolsep) * \real{0.2500}}
  >{\raggedright\arraybackslash}p{(\linewidth - 8\tabcolsep) * \real{0.1875}}
  >{\raggedright\arraybackslash}p{(\linewidth - 8\tabcolsep) * \real{0.1875}}@{}}
\toprule\noalign{}
\begin{minipage}[b]{\linewidth}\raggedright
Anchor
\end{minipage} & \begin{minipage}[b]{\linewidth}\raggedright
Symbol
\end{minipage} & \begin{minipage}[b]{\linewidth}\raggedleft
Value
\end{minipage} & \begin{minipage}[b]{\linewidth}\raggedright
Source
\end{minipage} & \begin{minipage}[b]{\linewidth}\raggedright
Used in
\end{minipage} \\
\midrule\noalign{}
\endhead
\bottomrule\noalign{}
\endlastfoot
Bosonic critical dimension & \(D_{\rm crit}\) & \(26\) & Polyakov action
(textbook) & \(26!\), \(26^3\), \(\mathrm{Li}_{26}\) \\
Pochhammer rising factorial & \((1)_{26}\) & \(26!\) &
\href{QCalcGeom.cpp\#L431}{QCalcGeom.cpp L431} & \(G\) closure
denominator \\
Newton derivative & \(\partial_r^{26}(1/r)\cdot r^{27}\) & \(26!\) &
\href{QCalcGeom.cpp\#L450}{QCalcGeom.cpp L450} & kinetic \& source
variation \\
Geometric folding & \((26!)^{-1/13}\) & sub-Planck scale &
\href{dpm_vacuum_manifold.py\#L1257}{dpm\_vacuum\_manifold.py L1257} &
radial moment factorization \\
Factorial barrier & \((26!)^2\) & \(1.626\times10^{53}\) & this paper &
hierarchy gap in \(G\) \\
\end{longtable}
}

\begin{center}\rule{0.5\linewidth}{0.5pt}\end{center}

\subsection{References}\label{references}

\begin{enumerate}
\def\labelenumi{\arabic{enumi}.}
\tightlist
\item
  Murphy, D. T., ``PAPER\_593: Newton's Constant via Void Coupling''
  (Star-Magic, Session 240).
\item
  Murphy, D. T., ``PAPER\_1159: Resonance Phase Closure
  \(\Phi_{\rm res}=5/6\)'' (Star-Magic, Session 246).
\item
  Murphy, D. T., ``PAPER\_1160: Time-Reversal Zone Closure
  \(F_{\rm TRZ}=1/10\)'' (Star-Magic, Session 247).
\item
  \href{QCalcGeom.cpp\#L308}{QCalcGeom.cpp L308-462} -- Pochhammer /
  26th radial derivative implementation (already canonical in codebase).
\item
  \href{dpm_vacuum_manifold.py\#L1251}{dpm\_vacuum\_manifold.py
  L1251-1258} -- \((26!)^{-1/13}\) geometric folding scale.
\item
  \href{AXIOMS_AND_THEOREMS.md\#L288}{AXIOMS\_AND\_THEOREMS.md L288} --
  \(26!\) factorial-barrier hierarchy claim.
\item
  \href{COMPLETE_UQFF_EQUATIONS_REFERENCE.md\#L1065}{COMPLETE\_UQFF\_EQUATIONS\_REFERENCE.md
  L1065-1082} -- VDS 26-rung ladder identification with bosonic critical
  dim.
\item
  M. Kaku, \emph{Strings, Conformal Fields, and M-Theory} (Springer, 2nd
  ed.) -- bosonic \(D_{\rm crit}=26\) derivation from no-ghost theorem.
\item
  \texttt{\_lagrangian\_rederivation\_outline.py} L154-156 (Star-Magic)
  -- original 22-torus hypothesis, now refuted.
\end{enumerate}

\begin{center}\rule{0.5\linewidth}{0.5pt}\end{center}

\textbf{Signed:} Daniel T. Murphy \textbf{Date:} May 10, 2026 (Session
248) \textbf{Repository:} Star-Magic, commit pending
\textbf{Verification:}
\((1)_{26} = \prod_{k=1}^{26} k = 403{,}291{,}461{,}126{,}605{,}635{,}584{,}000{,}000 = 26!\)
exact integer.

\end{document}
